\section{Transformaciones del modelo monolítico}
\label{sec:transformaciones}
\textit{Responsable: Fabricio.}

El modelo relacional único del proyecto de Base de Datos I se reparte en
tres sistemas de persistencia autónomos (MySQL, PostgreSQL y MongoDB), uno
por microservicio con base de datos. Ninguna tabla se comparte entre
motores. Las relaciones que en el monolito eran claves foráneas físicas se
transforman en \emph{referencias suaves} validadas por REST en tiempo de
escritura. Los enumerados que antes vivían como tipos \texttt{ENUM} en un
solo esquema se replican por servicio, sincronizados por
\texttt{docs/contratos/enums.md} para que los \texttt{JOIN} de Athena
funcionen.

\subsection{Implementación y decisiones}

\subsubsection{Descomposición por microservicio}
Las veintiuna relaciones del modelo original se distribuyen así:

\begin{itemize}
  \item \textbf{MS1 (MySQL)}: \texttt{persona}, \texttt{pasajero},
    \texttt{categoria\_migratoria}, \texttt{ticket}, \texttt{checkin},
    \texttt{equipaje}. Dueño de la identidad de los viajeros y de las
    transacciones comerciales.
  \item \textbf{MS2 (PostgreSQL)}: \texttt{vuelo}, \texttt{aerolinea},
    \texttt{aeronave}, \texttt{asiento}, \texttt{empleado},
    \texttt{tripulacion}, \texttt{operativo\_tierra},
    \texttt{opera\_tripulacion}. Dueño de la operativa aeroportuaria.
  \item \textbf{MS3 (MongoDB)}: colecciones \texttt{recursos},
    \texttt{incidencias} y \texttt{asignaciones}. Dueño de la
    infraestructura física del terminal y de la gestión de incidencias.
\end{itemize}

\subsubsection{Jerarquía \texttt{Persona} partida}
En el modelo E-R original, \texttt{Persona} es superclase de
\texttt{Pasajero} y también de \texttt{Personal} (a su vez superclase de
\texttt{Tripulacion} y \texttt{Operativo\_Tierra}). En una arquitectura de
microservicios no es viable compartir la tabla \texttt{persona} entre MS1
y MS2, pues cada servicio administra su propio motor. La transformación
adoptada es:

\begin{itemize}
  \item \textbf{MS1} conserva \texttt{persona} $\to$ \texttt{pasajero}
    (Class Table Inheritance). \texttt{Pasajero} tiene su propio
    \texttt{id\_persona} heredado y su clave alternativa por
    \texttt{(tipo\_documento, numero\_documento)}.
  \item \textbf{MS2} crea una tabla \texttt{empleado} con
    \texttt{nombre}, \texttt{apellido} y \texttt{fecha\_nacimiento}
    embebidos (equivalente a lo que Persona aportaba en el monolito), y
    sus dos subclases \texttt{tripulacion} (con \texttt{num\_licencia}) y
    \texttt{operativo\_tierra} (con \texttt{area\_operativa}). No hay
    tabla \texttt{persona} en MS2; los atributos de identidad se duplican
    de forma deliberada.
\end{itemize}

Esta duplicación es consciente y se acepta como el costo de la autonomía
por servicio. Los diagramas E-R publicados en \texttt{docs/er/} de cada
repositorio reflejan la duplicación.

\subsubsection{Referencias suaves entre servicios}
Los identificadores que en el monolito eran claves foráneas entre tablas
de distintos dominios se convierten en referencias validadas por REST:

\begin{itemize}
  \item \texttt{ms1.ticket.vuelo\_id} y \texttt{ms1.equipaje.vuelo\_id}
    apuntan a \texttt{ms2.vuelo.id}. En cada \texttt{POST /tickets} y
    \texttt{POST /equipajes}, MS1 llama a
    \texttt{GET /api/vuelos/\{id\}/exists} en MS2. Si el vuelo no existe
    o está \texttt{Cancelado}, MS1 responde \texttt{422} sin persistir.
  \item \texttt{ms3.incidencia.retrasa\_vuelos[].vuelo\_id} y
    \texttt{ms3.asignacion.vuelo\_id} apuntan a \texttt{ms2.vuelo.id}.
    MS3 valida los identificadores contra MS2 al crear la incidencia o
    la asignación.
  \item MS4 no persiste: agrega vía HTTP las respuestas de MS1, MS2 y
    MS3 al servir el manifiesto. Es el consumidor de todas las
    referencias suaves, no las valida ni las crea.
\end{itemize}

Los logs de los servicios registran cada llamada saliente
(\texttt{X-Request-Id} propagado por nginx), lo que constituye la
evidencia obligatoria de que un microservicio consume a otro.

\subsubsection{Enumerados replicados y contrato compartido}
El enunciado exige valores idénticos de enums entre los tres servicios
para que las agregaciones de Athena funcionen. Se centralizan en
\texttt{docs/contratos/enums.md} del repositorio de documentación:
\texttt{estado\_vuelo}, \texttt{tipo\_vuelo}, \texttt{tipo\_documento},
\texttt{alianza}, \texttt{estado\_boarding}, \texttt{estado\_acople},
\texttt{gravedad}, \texttt{tipo\_incidencia}, \texttt{tipo\_recurso},
\texttt{nombre\_categoria\_migratoria}, \texttt{area\_operativa},
\texttt{clase\_aeronave} (OACI) y \texttt{frecuencia\_radar}. Cada servicio
los implementa con la sintaxis nativa de su motor (PostgreSQL
\texttt{CREATE TYPE ... AS ENUM}, MySQL \texttt{ENUM} o \texttt{CHECK},
validación en la capa aplicativa de MongoDB), pero los \emph{strings}
coinciden byte a byte para que Athena pueda hacer \texttt{GROUP BY} y
\texttt{JOIN} sobre esas columnas sin normalizaciones intermedias. Los
valores no llevan tildes en enums (\texttt{Critica}, \texttt{Transito},
\texttt{Inundacion}) para evitar problemas de codificación en la ruta
CSV $\to$ Glue $\to$ Athena.

\subsubsection{Rangos de identificadores deterministas}
Para que los \texttt{JOIN} entre tablas de tres motores distintos
funcionen en Athena, los identificadores deben cruzarse aunque la
generación de datos ocurra en tres corridas independientes por servicio.
Se acuerdan en \texttt{docs/plan-de-trabajo.md \S3} rangos disjuntos por
entidad (\texttt{vuelo.id} 1--25\,000, \texttt{persona.id} 100\,000--160\,000,
\texttt{ticket.id} 1--60\,000, \texttt{recurso.id} 1--500,
\texttt{incidencia.id} 1--30\,000) y un \texttt{SEED} fijo
(\texttt{20\,260\,905}) para el generador \texttt{seeds/} compartido. El
orden interno del generador es \textbf{MS2 $\to$ MS1 $\to$ MS3} para que
los identificadores de vuelo existan antes de que se creen los tickets y
las incidencias que los referencian. Detalles en la Sección
\ref{sec:data-science}.

\subsection{Resultados y evidencias}

Las transformaciones descritas se materializan en artefactos verificables:

\begin{itemize}
  \item Los tres diagramas E-R por servicio residen en \texttt{docs/er/}
    del repositorio de documentación. Cada uno refleja únicamente las
    tablas de su motor, con la duplicación deliberada de atributos de
    identidad.
  \item Los archivos de contrato compartido: \texttt{docs/contratos/enums.md}
    (enumerados) y \texttt{docs/plan-de-trabajo.md \S3} (rangos de ID).
  \item Los generadores en \texttt{aeropuerto-data-science/seeds/}
    respetan los rangos y el orden acordado, y producen salida
    reproducible byte a byte con el mismo \texttt{SEED} (comprobado por
    igualdad de sumas MD5 entre corridas).
  \item La validación cruzada de las referencias suaves se ejecuta con
    las cinco consultas Q1--Q5 sobre Postgres local (Sección
    \ref{sec:data-science}); si un identificador no cruzara, los
    \texttt{JOIN} devolverían cero filas. Los resultados no vacíos
    confirman la integridad.
\end{itemize}
